% Part: incompleteness
% Chapter: representability-in-q
% Section: comp-representable

\documentclass[../../../include/open-logic-section]{subfiles}

\begin{document}

\olfileid[id]{inc}{req}{crq}
\olsection{Fungsi yang Dapat Dihitung Dapat Direpresentasikan dalam $\Th{Q}$}

\begin{thm}
Setiap fungsi yang dapat dihitung dapat direpresentasikan dalam~$\Th{Q}$.
\end{thm}

\begin{proof}
Agar pasti, dan dengan menggunakan Tesis Church--Turing, katakanlah bahwa
suatu fungsi dapat dihitung jika dan hanya jika fungsi itu rekursif umum.
Fungsi-fungsi rekursif umum ialah fungsi-fungsi yang dapat didefinisikan dari
fungsi nol~$\Zero$, fungsi penerus~$\Succ$, dan fungsi proyeksi~$\Proj{n}{i}$
dengan menggunakan komposisi, rekursi primitif, dan pencarian tak berbatas
pada fungsi reguler. Menurut \olref[pri]{lem:prim-rec}, setiap fungsi~$h$
yang dapat didefinisikan dari $f$ dan~$g$ dengan rekursi primitif juga dapat
didefinisikan melalui komposisi dan pencarian tak berbatas pada fungsi
reguler dari $f$, $g$, dan $\Zero$, $\Succ$, $\Proj{n}{i}$, $\Add$, $\Mult$,
$\Char{=}$. Akibatnya, suatu fungsi rekursif umum jika dan hanya jika fungsi
itu dapat didefinisikan dari $\Zero$, $\Succ$, $\Proj{n}{i}$, $\Add$,
$\Mult$, $\Char{=}$ dengan menggunakan komposisi dan pencarian tak berbatas
pada fungsi reguler.

Selain itu, telah kita tunjukkan bahwa fungsi-fungsi dasar yang bersangkutan
dapat direpresentasikan dalam~$\Th{Q}$
(\cref{inc:req:bre:prop:rep-zero,inc:req:bre:prop:rep-succ,inc:req:bre:prop:rep-proj,inc:req:bre:prop:rep-id,inc:req:bre:prop:rep-add,inc:req:bre:prop:rep-mult}),
dan bahwa setiap fungsi yang didefinisikan dari fungsi-fungsi yang dapat
direpresentasikan melalui komposisi atau pencarian tak berbatas pada fungsi
reguler (\olref[cmp]{prop:rep-composition},
\olref[min]{prop:rep-minimization}) juga dapat direpresentasikan. Jadi, setiap
fungsi rekursif umum dapat direpresentasikan dalam~$\Th{Q}$.
\end{proof}

\begin{explain}
Telah kita tunjukkan bahwa himpunan fungsi yang dapat dihitung dapat
dicirikan sebagai himpunan fungsi yang dapat direpresentasikan dalam
$\Th{Q}$. Sebenarnya, pembuktian tersebut lebih umum. Dari definisi
representasi, tidak sulit melihat bahwa setiap teori yang memperluas
$\Th{Q}$ (atau yang memungkinkan $\Th{Q}$ ditafsirkan di dalamnya) dapat
merepresentasikan fungsi-fungsi yang dapat dihitung. Namun, sebaliknya, dalam
setiap sistem !!{derivation} yang gagasan !!{derivation}-nya dapat dihitung,
setiap fungsi yang dapat direpresentasikan dapat dihitung. Jadi, misalnya,
himpunan fungsi yang dapat dihitung dapat dicirikan sebagai himpunan fungsi
yang dapat direpresentasikan dalam aritmetika Peano, atau bahkan teori
himpunan Zermelo--Fraenkel. Seperti dicatat G\"odel, hal ini agak
mengejutkan. Akan kita lihat bahwa dalam perkara keterbuktian, jawabannya
sangat peka terhadap teori mana yang dipertimbangkan; secara kasar, semakin
kuat aksiomanya, semakin banyak yang dapat dibuktikan. Namun, dalam cakupan
luas teori-teori aksiomatik, fungsi-fungsi yang dapat direpresentasikan tepat
merupakan fungsi-fungsi yang dapat dihitung; teori yang lebih kuat tidak
merepresentasikan lebih banyak fungsi selama teori tersebut dapat
diaksiomatisasi.
\end{explain}

\end{document}
